% Copyright (c) 2026 Eugene Joseph M. Ragasa.
% SPDX-License-Identifier: Apache-2.0
% Author: Eugene Joseph M. Ragasa
\documentclass[11pt]{article}
\usepackage[T1]{fontenc}
\usepackage{amsmath,amssymb}
\usepackage{booktabs,tabularx}
\usepackage{seqsplit}
\usepackage[margin=1in]{geometry}
\usepackage[hidelinks]{hyperref}
\usepackage{listings}
\newcolumntype{Y}{>{\raggedright\arraybackslash}X}
\lstset{basicstyle=\ttfamily\small,breaklines=true,columns=fullflexible,literate={•}{{$\mathbin{\bullet}$}}1}
\title{Project Koios Formal Proof Dossier\\\large PIAB1D-SPECTRUM-THM-015: canonical finite Dirichlet mode eigenpair}
\author{Eugene Joseph M. Ragasa}
\date{Formal proof completed: 2026-10-21}
\begin{document}
\maketitle

\section{Identity and status}
\begin{tabularx}{\textwidth}{@{}lY@{}}\toprule
Statement ID & \texttt{PIAB1D-SPECTRUM-THM-015} \\
Status & proved; no admissions \\
Lean declaration & \path{ProjectKoios.QuantumMechanics.Piab1d.negativeDirichletDifferenceMatrix_mulVec_canonicalMode_cosine} \\
Owning module & \path{ProjectKoios/QuantumMechanics/Piab1d/Spectrum.lean} \\
Author and license & Eugene Joseph M. Ragasa; Apache-2.0 \\
Formal proof completed & 2026-10-21 \\
\bottomrule\end{tabularx}
Status concerns only the exact Lean proposition. It does not imply numerical validation,
source fidelity, novelty, applicability, publication, or release.

\section{Exact mathematical statement}
For natural numbers $N$ and $k$, let
\[
 z_{N,k}=\exp\!\left(i\frac{k\pi}{N+1}\right),\qquad
 v_j=\frac{z_{N,k}^{j+1}-z_{N,k}^{-(j+1)}}{2i},\quad 0\leq j<N.
\]
For the exact finite negative centered-difference matrix $D_N$ with zero Dirichlet ghost
values,
\[
 D_Nv=\left(2-2\cos\frac{k\pi}{N+1}\right)v.
\]
The Lean theorem allows every natural $k$ and therefore does not assert that every such vector
is nonzero or that all indices give distinct modes.

\section{Source provenance and authority boundary}
\begin{tabularx}{\textwidth}{@{}lY@{}}\toprule
Owner repository & \texttt{ksdft2effmass} \\
Navigation branch & \texttt{work/source-package-ownership-cleanup} \\
Base commit & \texttt{2229777163882c3452b950b8b21db86c7ea43511} \\
Dirty-state limitation & the PIAB1D ownership restructuring was uncommitted at inspection \\
Exact inspected-byte hashes & recorded in \path{docs/source-audit.md} \\
\bottomrule\end{tabularx}
The theorem is source-neutral exact mathematics. It neither authenticates owner execution nor
imports binary64 arrays or retained JSON as formal evidence.

\section{Formulation-specific possible prior art}
The result and the central proof formulation are classical and carry no novelty claim. Our Lean
proof should be understood as a formal synthesis of two established viewpoints: solve the centered
recurrence using counter-propagating complex characteristic modes, and obtain homogeneous
Dirichlet modes as the antisymmetric or odd sector of a periodic Fourier representation. The
following inspected works are the closest formulation-level precedents located so far.

\begin{itemize}
\item \textbf{Yueh (2005), complex characteristic-mode derivation.} Wen-Chyuan Yueh,
``Eigenvalues of Several Tridiagonal Matrices,'' \emph{Applied Mathematics E-Notes} 5
(2005), 66--74. Section~2, pp.~67--70, embeds the finite eigenvector into a complex
recurrence with ghost values $u_0=u_{n+1}=0$, writes the two characteristic roots as
$\gamma_{\pm}=\rho^{-1}e^{\pm i\theta}$, subtracts their powers, and applies De Moivre's
theorem to obtain sine modes. Section~3, especially equations (11)--(12) on pp.~70--71,
imposes $\sin((n+1)\theta)=0$ and gives the corresponding tridiagonal eigenvectors and
eigenvalues. This closely precedes the $z^n-z^{-n}$ portion of our proof, although it does
not use our source-neutral Lean decomposition into translation, periodicity, reality, and
\texttt{Matrix.mulVec} bridge lemmas. Inspected \href{https://www.math.nthu.edu.tw/~amen/2005/040903-7.pdf}{candidate PDF};
SHA-256 \texttt{\seqsplit{f2a2f513fc9cc362bf47831bc1696b6f6ea54651d74dc0e04b9c47001b4de4c6}}.

\item \textbf{Püschel--Moura (2008), antisymmetric Fourier/DST viewpoint.} Markus Püschel and
José M. F. Moura, ``Algebraic Signal Processing Theory: 1-D Space,'' \emph{IEEE
Transactions on Signal Processing} 56(8) (2008), 3586--3599. Sections~III--IV,
pp.~3590--3593, treat Dirichlet ``zero value'' boundaries through whole- or half-point
antisymmetric extensions and derive finite discrete trigonometric transforms as Fourier
transforms of the associated finite space models. This is formulation-level prior art for
our decision to derive Dirichlet sine modes from an antisymmetric sector rather than treat
sine as an unrelated primitive. Inspected
\href{https://users.ece.cmu.edu/~moura/papers/t-sp-aug08-part2--pueschelmoura-ieeexplore.pdf}{candidate PDF};
SHA-256 \texttt{\seqsplit{a3132903efc58592fbbb795da982adc2fc561f203582aa22f066591fef79199e}}.

\item \textbf{Noschese--Pasquini--Reichel (2013), broader spectral context.} Silvia Noschese,
Lionello Pasquini, and Lothar Reichel, ``Tridiagonal Toeplitz Matrices: Properties and
Novel Applications,'' \emph{Numerical Linear Algebra with Applications} 20 (2013),
display DOI \url{https://doi.org/10.1002/nla.1811}. This is retained only as a broader
candidate survey stating that the closed-form tridiagonal Toeplitz spectrum is known; it
has not yet been inspected as evidence for our particular exponential/antisymmetric proof
architecture.
\end{itemize}

No source has yet been located and audited that packages every step exactly as this Lean tranche
does: a reusable integer-lattice character $z^n$, root-of-unity periodicity, the antisymmetric
combination $z^n-z^{-n}$, explicit real representability after division by $2i$, an exact finite
\texttt{Matrix.mulVec} realization, and the scaling bridge to the declared PIAB1D Hamiltonian.
That packaging distinction is not a mathematical novelty claim. These candidates were not used as
axioms or external proof steps. Their References-owned canonical identities remain unresolved, so
a completed literature crosswalk still requires identity resolution and an audited
statement/assumption comparison.

\section{Proof schematic}
The reusable complex-exponential layer first proves the translation recurrence
\[
-(z^{n-1}-2z^n+z^{n+1})=(2-z^{-1}-z)z^n.
\]
Linearity transfers it to $z^n-z^{-n}$ and division by $2i$. The phase is nonzero, is a
$2(N+1)$-st root of unity, and hence makes the antisymmetric mode vanish at both ghost
endpoints. Restriction to interior coordinates and the exact \texttt{toMatrix'}--\texttt{mulVec}
bridge give the finite matrix equation. Finally,
$z+z^{-1}=2\cos(k\pi/(N+1))$ rewrites the eigenvalue.

\section{Lean realization}
\begin{lstlisting}
theorem negativeDirichletDifferenceMatrix_mulVec_canonicalMode_cosine
    (N k : Nat) :
    Matrix.mulVec (negativeDirichletDifferenceMatrix N)
        (normalizedDirichletExponentialModeVector N (dirichletPhase N k)) =
      ((2 - 2 * Real.cos (dirichletPhaseAngle N k) : Real) : Complex) •
        normalizedDirichletExponentialModeVector N (dirichletPhase N k)
\end{lstlisting}
The repository source is authoritative; this is a shortened rendering.

\section{Toolchain and trusted-base report}
\begin{tabularx}{\textwidth}{@{}lY@{}}\toprule
Lean toolchain & Lean 4.35.0 \\
Mathlib revision & \texttt{56a005855ea17e0b0958f73ec801ae140b0a1679} \\
Build & \texttt{lake build}: passed, 3,124 jobs \\
Admission scan & \texttt{./scripts/audit.sh}: passed \\
Axioms & \texttt{propext}, \texttt{Classical.choice}, \texttt{Quot.sound} \\
\bottomrule\end{tabularx}
No \texttt{sorry}, \texttt{admit}, or repository-owned scientific axiom occurs.

\section{AI, automation, and review}
Proof discovery, Lean implementation, and dossier drafting were AI-assisted. Lean's elaborator,
kernel, Mathlib tactics, Lake, and repository scripts checked the exact declaration and reported
axioms. Eugene Joseph M. Ragasa selected the controlled-model question, operated the checks, and
reviewed the stated authority boundary. No independent expert review is recorded.

\section{Explicit nonclaims}
This theorem does not prove continuum convergence, completeness or orthogonality of a spectral
basis, physical parameter applicability, correspondence to a NumPy/SciPy matrix, correctness of
floating-point eigenpairs, retained-result validity, or Bloch--Floquet theory. The common
complex-exponential representation is reusable infrastructure for a later separately reviewed
Bloch bridge, not that bridge itself.

\section{Revision history}
\begin{tabularx}{\textwidth}{@{}lllY@{}}\toprule
Date & Revision & Process & Outcome \\
\midrule
2026-10-21 & initial & Lean kernel and repository audit & exact proposition checked; no independent expert review \\
\bottomrule\end{tabularx}
\end{document}
